\section{Discussion and Limitations}
\label{sec:discussion}

\subsection{Design reflections}
Three choices paid off. Making the causaloid the per-series unit kept the detector
small and the control flow legible, since a detection reads as a pipeline rather
than nested conditionals. Splitting a reusable core from the per-series engine gave
verdict parity for free: the same code computes the score wherever it runs, so we
test the core once and trust it everywhere. Persisting only the window, and
rebuilding the statistics each pass, made per-series state trivial to checkpoint and
removed any question of a stored summary disagreeing with the samples it
represents. The price is a recompute over the window, and \cref{sec:eval} shows
that price is small enough to ignore at the rates an agent produces.

The deeper point is that a cheap reasoner changes where detection can live. When
scoring a sample costs a fraction of a microsecond, detection is a rounding error in
the host's budget, so it can run on every agent instead of in a central tier that
grows with the fleet.

\subsection{Limitations}
The honest gap is accuracy. Every number in this paper is throughput or memory. We
have not measured precision or recall, because we do not yet have labeled anomalies
to measure against, so we cannot claim the detector catches the right things, only
that it runs cheaply. The benchmark workload is synthetic, taken on a single
machine in a single run, and the resident-size figures come from a platform that is
not the deployment target. The method itself is modest by design: a $z$-score over a
few hundred samples assumes the series is roughly stationary over that window and
finds point and level shifts rather than slow seasonal structure. That is the
intent of the split: the edge engine handles fast, short-window anomalies, while
seasonal and longer-horizon detection is the job of the central tier in
\cref{sec:seasonal}, which composes with the edge as recall and precision.
Finally, the engine detects but does not act: the corrective step that \dc{}'s
example takes is deliberately left off.

\subsection{Future work}
The most valuable next step is a detection-quality study: inject labeled anomalies,
measure precision and recall, and compare against at least an EWMA baseline and one
seasonal method. A reasoner-only microbenchmark would separate the causaloid's cost
from the per-series engine, and the reversible $O(1)$ accumulator in the core is the
obvious lever if the recompute ever becomes the bottleneck. Beyond that, the
framework's intervention step opens the door to guarded corrective actions on the
host.
